%! Author = Joey Dufourd
%! Date = 30/10/2025
\documentclass[11pt,a4paper]{article}
\usepackage[margin=2.5cm]{geometry}
\usepackage[T1]{fontenc}
\usepackage[utf8]{inputenc}
\usepackage{lmodern}
\usepackage{setspace}
\usepackage[hidelinks]{hyperref}
\usepackage{parskip}
\setstretch{1.15}

%pdflatex -output-directory=assets MotivationLetter.tex

\begin{document}

\begin{flushright}
Joey Dufourd, Ph.D.\\
Institut de Génétique Humaine (CNRS)\\
Montpellier, France\\
\href{mailto:joey.dufourd@gmail.com}{joey.dufourd@gmail.com}\\
\href{https://joeydufourd.github.io}{https://joeydufourd.github.io}\\
\vspace{0.4cm}
\today
\end{flushright}

\vspace{0.4cm}
\noindent
\textbf{Subject: Application for a Postdoctoral Research Position in Computational Biology}

\vspace{0.5cm}
Dear [Dr./Professor/Selection Committee],

\vspace{0.3cm}
I am writing to express my sincere interest in joining your laboratory as a postdoctoral researcher in computational biology. I have recently completed my Ph.D. at the Institut de Génétique Humaine (CNRS, Montpellier), where I studied the molecular mechanisms that safeguard telomeres and influence chromatin organization. My work brought together molecular biology, super-resolution imaging, and computational analysis, allowing me to explore genome regulation from both experimental and data-driven perspectives.

During my doctoral work, I identified and characterized \textit{TELS1}, a pluripotency-specific telomere-binding protein that stabilizes \textit{t}-loops independently of TRF2. This finding points toward a previously unrecognized mechanism of telomere protection in early development. To better understand this process, I built analysis pipelines using Python and R, including AI-based tools to automate 3D-SIM image quantification and integrate imaging, proteomic, and molecular datasets. These projects gave me solid experience in image analysis, data integration, and machine learning, while strengthening my belief that computational tools can reveal hidden biological principles.

Alongside my research, I have been deeply involved in community and teaching activities. I co-founded and coordinate the IGH Young Scientist Club to encourage collaboration and exchange between early-career researchers. I have also taught courses in genetics and epigenetics, as well as workshops introducing AI and statistics to biologists. Sharing science and helping others grow as researchers is something I truly value.

What motivates me now is to apply these experiences to larger-scale, computational studies of genome regulation. Your group’s work on [specific topic or project] strongly resonates with my own interests. I am particularly eager to combine my background in telomere biology with advanced data analysis to explore new aspects of chromatin dynamics and cell identity. I enjoy working at the interface between wet-lab biology and computation—where new questions often arise.

I am a curious, focused, and enthusiastic scientist, and I approach my work with both rigor and creativity. I would be delighted to contribute to your research environment and to learn from your team’s expertise. My CV is attached, and I would be glad to provide any additional information or references if needed.

\vspace{0.4cm}
\noindent
Yours sincerely,\\[0.5cm]
\textbf{Joey Dufourd}

\end{document}
